\section{Contoh Program}
\label{sec:contoh-bab3}

Bagian ini menyajikan tiga contoh program yang dirakit dan dijalankan
menggunakan Emu8086 atau TASM. Ketiganya sengaja dibuat sesingkat
mungkin agar seluruh baris dapat dipahami satu per satu.

\subsection{Contoh 1: Mencetak Teks dengan Layanan DOS}

Program pada Kode~\ref{lst:hello} adalah program assembly pertama yang
biasa dipakai: mencetak satu baris teks ke layar.

\lstinputlisting[
    language=MASM,
    caption={Program \file{hello.asm} mencetak teks ke layar.},
    label={lst:hello}
]{\CodeDir/hello.asm}

Penjelasan setiap bagian program:

\begin{itemize}
    \item \lstasm|.MODEL SMALL| memberi tahu assembler bahwa program
          memakai model memori kecil.
    \item \lstasm|.DATA| membuka segmen data; di dalamnya terdapat
          variabel \code{pesan} berupa deretan byte yang diakhiri
          \code{0Dh}, \code{0Ah} (pindah baris), dan \code{'\$'}.
    \item \lstasm|panjang EQU \$ - pesan| menghitung panjang string
          secara otomatis dari selisih posisi sekarang dengan awal
          label \code{pesan}.
    \item \lstasm|MOV AH, 09h| dan \lstasm|LEA DX, pesan| menyiapkan
          argumen layanan DOS: \reg{AH} berisi nomor fungsi, \reg{DX}
          berisi alamat string.
    \item \lstasm|INT 21h| menyerahkan kendali ke layanan DOS untuk
          mencetak string.
    \item \lstasm|MOV AH, 4Ch| dan \lstasm|INT 21h| mengakhiri program.
\end{itemize}

\begin{session}[caption={Keluaran program \file{hello.asm}.}, label={lst:keluaran-hello}]
Halo, dunia!
\end{session}

\subsection{Contoh 2: Program \file{.COM} dengan \dir{ORG 100h}}

Bentuk lain yang lebih ringkas adalah program \file{.COM}. Karena
seluruh kode, data, dan stack berada dalam satu segmen, inisialisasi
\reg{DS} tidak diperlukan. Program ditunjukkan pada Kode~\ref{lst:com-org}.

\begin{lstlisting}[
    language=MASM,
    caption={Program \file{.COM} ringkas dengan \dir{ORG 100h}.},
    label={lst:com-org}
]
.MODEL TINY
.CODE
ORG 100h            ; program .COM dimulai pada offset 100h

START:
    MOV AH, 09h     ; cetak string
    MOV DX, OFFSET pesan
    INT 21h

    MOV AH, 4Ch     ; keluar
    MOV AL, 0
    INT 21h

pesan DB 'Selamat datang di assembly!', 0Dh, 0Ah, '$'
END START
\end{lstlisting}

\begin{infobox}
Perhatikan bahwa pada program \file{.COM} tidak ada
\lstasm|MOV AX, @DATA| maupun \lstasm|MOV DS, AX|. DOS sudah
menetapkan \reg{DS} ke segmen program, sehingga variabel dapat
langsung dialamatkan.
\end{infobox}

\subsection{Contoh 3: Menampilkan Karakter Satu per Satu}

Layanan DOS fungsi \hex{02} mencetak satu karakter yang berada pada
register \reg{DL}. Contoh ini menunjukkan bagaimana sebuah huruf
dapat dicetak, kemudian diubah, lalu dicetak kembali
(Kode~\ref{lst:karakter}).

\begin{lstlisting}[
    language=MASM,
    caption={Mencetak karakter dengan \code{INT 21h} fungsi \hex{02}.},
    label={lst:karakter}
]
.MODEL SMALL
.STACK 100h

.CODE
START:
    MOV AX, @DATA
    MOV DS, AX

    MOV AH, 02h     ; fungsi 02h : cetak karakter pada DL
    MOV DL, 'A'
    INT 21h

    MOV DL, 'B'
    INT 21h

    MOV DL, 'C'
    INT 21h

    MOV AH, 4Ch
    MOV AL, 0
    INT 21h

END START
\end{lstlisting}

\begin{session}[caption={Keluaran program pada Kode~\ref{lst:karakter}.}]
ABC
\end{session}

Ketiga contoh di atas menunjukkan pola yang sama: siapkan argumen
pada register, panggil layanan melalui \code{INT 21h}, lalu akhiri
program. Pola inilah yang akan dipakai ulang pada bab-bab berikutnya.
